\section{Schlussfolgerungen und Ausblick}\label{sec:schluss}

\subsection{Fazit}\label{sec:fazit}

Diese Arbeit zeigt, dass ein kleines Modell die oberste Karte eines
Jassstapels auf dem Smartphone in Echtzeit erkennen kann. Das Modell
wurde mit synthetischen Bildern aus den Scans zweier Kartenblätter
trainiert. In der letzten Datenmengenreihe erreicht der einstufige
Detektor C bei 200'000 Trainingsbildern auf dem realen
Validierungsset eine \map{50} von 0.908. Seine Leistung steigt bis
dorthin mit jeder Vergrösserung des Datensatzes. Ein zusätzlicher Lauf
mit 500'000 Bildern brachte keinen weiteren Gewinn. Die Menge der Bilder
aus diesem Generator ist damit ausgereizt.

Für die App ist C die am besten belegte Wahl. Es erkennt Klasse und
Position in einem Schritt. Die reine Klassifikation A liefert keine
Kartenposition. Für die vollständige zweistufige Erkennung B und
einen gemeinsamen Latenzvergleich liegen keine verlässlichen
Messungen aus der letzten Reihe vor. Deshalb lässt sich kein genauer
Leistungsabstand zwischen C und B angeben.

Der Transfer von synthetischen Trainingsbildern auf reale Aufnahmen
funktioniert auf dem vorhandenen Validierungsset. Wie zuverlässig
die App einen ganzen Stapel in neuen Spielsituationen zählt, ist
damit noch nicht gezeigt. Neben dem Modell sind ein Datengenerator,
eine reale Prüfbasis, eine Trainingspipeline und Zähl-Apps für iOS
und Android entstanden. Die Apps wurden in den Stores eingereicht.

\subsection{Ausblick}\label{sec:ausblick}

Als Nächstes soll das System mit neuen, unabhängigen Aufnahmesessions
geprüft werden. Dafür sollen ganze Zählvorgänge direkt in der App als
Video aufgezeichnet werden. Das Dataset Tool kann diese Sessions bereits
erneut abspielen und auswerten (Abbildung~\ref{fig:analyzer}). Künftig
soll diese Auswertung stärker automatisiert werden. Auffällige Frames
können direkt ins Validierungsset übernommen werden. Dadurch wird
sichtbar, welche Situationen in den Trainingsdaten noch fehlen und
gezielt ergänzt werden müssen.

\begin{figure}[H]
\centering
\includegraphics[width=\textwidth]{analyzer.png}
\caption{Auswertung einer aufgezeichneten Zählsession im Dataset Tool. Links
ist die Vorhersage für das ausgewählte Frame zu sehen. Rechts werden
auffällige Erkennungen und häufige Verwechslungen zusammengefasst. Über die
Zeitachse am unteren Rand lässt sich die gesamte Aufnahme prüfen. Das
ausgewählte Frame zeigt dieselbe falsche Erkennung wie das linke Bild in
Abbildung~\ref{fig:frames-nacheinander} und kann direkt ins Validierungsset
übernommen werden.}
\label{fig:analyzer}
\end{figure}

Auch die Labels des Validierungssets müssen möglicherweise genauer
unterteilt werden. Beim Labeln stellte sich immer wieder die Frage, ob
eine schwierige Karte noch erkannt werden muss oder bereits als
unlesbar gilt. Diese Entscheidung beeinflusst die Auswertung direkt.
Bei einigen Grenzfällen sind sowohl eine Erkennung als auch keine
Erkennung nachvollziehbar. Solche Aufnahmen könnten künftig gesondert
markiert und getrennt ausgewertet werden.

Für solche Grenzfälle gibt es ein etabliertes Vorbild. Die
PASCAL-VOC-Challenge markiert Objekte, die wegen ihrer Grösse, der
Beleuchtung, der Bildqualität oder des nötigen Kontextwissens besonders
schwer zu erkennen sind, mit einem eigenen Attribut. In der Auswertung
werden diese Objekte verworfen, und wer sie dennoch erkennt, wird dafür
nicht bestraft \cite{everingham2010pascal}.

YOLO kennt ein solches Attribut nicht. Das Labelformat sieht es nicht vor,
und Ultralytics verweist bei dieser Frage auf einen Eingriff in die
Verlustfunktion \cite{ultralytics2025ignore}. Die Markierung müsste
deshalb neben dem Validierungsset geführt und in einer eigenen Auswertung
berücksichtigt werden.

Aktuell werden beide Kartenblätter mit demselben Modell erkannt. Für
weitere Kartensets könnten stattdessen eigene Modelle pro Blatt
trainiert werden. Die App kennt das gewählte Kartenblatt bereits und
könnte das passende Modell laden. Dadurch liessen sich zusätzliche
Kartensets unabhängig voneinander ergänzen, ohne dass jedes Modell
immer mehr Klassen lernen muss.

Die Beschränkung auf die oberste Karte stellt zudem Anforderungen an
das Ablegen der Karten. Eine mögliche Erweiterung wäre, mehrere
vollständig sichtbare Karten gleichzeitig zu erkennen. YOLO kann
mehrere Objekte in einem Bild ausgeben. Die App müsste dabei
sicherstellen, dass jede Karte nur einmal gezählt wird.

Beim Testen zeigte sich ausserdem, dass das gleichzeitige Halten des
Smartphones und Ablegen der Karten umständlich sein kann. Zwei Personen
können diese Aufgabe aufteilen. Für eine einzelne Person wäre eine
einfache Halterung hilfreich. Sie könnte das Smartphone im richtigen
Winkel über dem Tisch halten und darunter eine Ablage für die gezählten
Karten bieten. Eine solche Vorrichtung würde die Bedienung vereinfachen
und für gleichmässigere Aufnahmen sorgen.
